\documentclass{article}
%%%%% PROGRAMMING PACKAGES
\usepackage{expl3,xparse}     %LaTeX3
\usepackage{mfirstuc}    %first letter uppercase with \makefirstuc (used in todo notes)
\usepackage{settobox}     %get box dimensions with \settoboxwidth (used in fit environment) => pretty stupid package.. could be avoided

%%%%% FIGURES, BOXES, COLORS
\usepackage{graphics,graphicx,trimclip,adjustbox}     %figures and box management
\usepackage{color,xcolor}     %to use colors (used in the todonotes)
    
%%%%% ENUMERATE
\usepackage{enumitem}   %to use "enumerate"
    \setlist[enumerate,1]{label = $\bullet$}   %default label at bullet points
    
%%%%% TODO NOTES
\usepackage[textwidth=16.3mm,textsize=scriptsize]{todonotes}   %to put notes everywhere
\newcommand*{\newcommentator}[2][red!20]{%
    \expandafter\newcommand\csname #2\endcsname[1]{\todo[color=#1]{{\bf\makefirstuc{#2}:} ##1}}%
    \expandafter\newcommand\csname big#2\endcsname[1]{\todo[inline,color=#1]{{\bf\makefirstuc{#2}:} ##1}}}
\newcommentator{pablo}

%%%%% MY HYPERLINK AND REF/LABEL MACROS
\usepackage{mypreamble/linkref}

%%%%% MATH PACKAGES
\usepackage{amssymb,amsmath}     %math original package
    \numberwithin{equation}{section}
\usepackage{mathtools}     %improves AMSmath
\usepackage{stackrel}     %better stacking of math symbols
\usepackage{mleftright}     %defines \mleft and \mright
    \mleftright     %redefines \left as \mleft and \right as \mright
\usepackage{pgfcore}
\usepackage{tikz-cd}     %advanced and rigorous diagrams
    
%%%%% MY SPECIAL PACKAGES
\usepackage{mypreamble/mathenv}     %improves on AMSmath's equation
\usepackage{mypreamble/thmenv}     %replaces the AMSthm package

%%%%% MATH MACROS AND FONT PACKAGES
%\usepackage{fontspec}
\usepackage{mypreamble/mathcro}

%%%%% HYPERREF PACKAGE
\makeatletter
\pretocmd\refstepcounter{%
 \zref@setcurrent{counter}{#1}}{}{}
\makeatother
\usepackage[pdfencoding=auto]{hyperref}   %hyperlinks parts of the pdf
    \hypersetup{
        bookmarksopen=true,   %expends tree of bookmarks in Acrobat pdf reader
        bookmarksnumbered=true,  %show numbers of sections in bookmarks
        linktoc=all,   %both section and page numbers are links
        hidelinks   %links (color+border) are hidden
    }
\providecommand*{\definitionautorefname}{Definition}
\providecommand*{\notationautorefname}{Notation}
\providecommand*{\conventionautorefname}{Convention}
\providecommand*{\conjectureautorefname}{Conjecture}
\providecommand*{\assumptionautorefname}{Assumption}
\providecommand*{\remarkautorefname}{Remark}
\providecommand*{\observationautorefname}{Observation}
\providecommand*{\lemmaautorefname}{Lemma}
\providecommand*{\propositionautorefname}{Proposition}
\providecommand*{\corollaryautorefname}{Corollary}
\providecommand*{\exampleautorefname}{Example}
\providecommand*{\exerciseautorefname}{Exercise}
\providecommand*{\openquestionautorefname}{Open Question}

%%%%% OUTPUT LAYOUT PACKAGES
\usepackage{geometry}     %to set margins
\geometry{
    a4paper,
    total={170mm,257mm},
    left=20mm,
    top=20mm,
    marginparwidth=20mm,
}
%% Clearpage before sections but not before TOC
\let\oldsection\section
\renewcommand\section{\clearpage\oldsection}
\makeatletter
\renewcommand\tableofcontents{%
    \oldsection*{\contentsname\@mkboth{\MakeUppercase\contentsname}{\MakeUppercase\contentsname}}\@starttoc{toc}%
}
\makeatother    %ensures the TOC stays in front page

%%%%% TYPESET DEBUGGING
%\usepackage{showframe}     %to make margins appear
%\usepackage{lua-visual-debug}     %to make all boxes appear
%\overfullrule=1pt   %creates 1pt wide black bars next to any overfull hbox

%%%%% QUIVER SETUP
\usetikzlibrary{calc,decorations.pathmorphing}
% A TikZ style for curved arrows of a fixed height, due to AndréC.
\tikzset{curve/.style={settings={#1},to path={(\tikztostart)
    .. controls ($(\tikztostart)!\pv{pos}!(\tikztotarget)!\pv{height}!270:(\tikztotarget)$)
    and ($(\tikztostart)!1-\pv{pos}!(\tikztotarget)!\pv{height}!270:(\tikztotarget)$)
    .. (\tikztotarget)\tikztonodes}},
    settings/.code={\tikzset{quiver/.cd,#1}
        \def\pv##1{\pgfkeysvalueof{/tikz/quiver/##1}}},
    quiver/.cd,pos/.initial=0.35,height/.initial=0}

% TikZ arrowhead/tail styles.
\tikzset{tail reversed/.code={\pgfsetarrowsstart{tikzcd to}}}
\tikzset{2tail/.code={\pgfsetarrowsstart{Implies[reversed]}}}
\tikzset{2tail reversed/.code={\pgfsetarrowsstart{Implies}}}
% TikZ arrow styles.
\tikzset{no body/.style={/tikz/dash pattern=on 0 off 1mm}}
\usepackage{stmaryrd}

%%%%% TITLE PAGE
\title{Action on Categories with Adjoints}
\author{Pablo Bustillo Vazquez}

\begin{document}

\maketitle
\tableofcontents
\newpage

\section{Ideas}

I'm starting to grow convinced that the action will come from applying the "classifying space" functor to a map of posets.
More specifically, there's many colimit diagrams describing $c_k^{n-\adj}$ and they are arranged in a poset.
I think roughly moving through that poset means describing the same cell by different formulae.
This should correspond to the automorphisms of categories with adjoints.

If my "signed Schubert cells" are correct, then $BO(n)$ is the realization of a poset.
A reasonnable question is: can I find a similar poset for $PL(n)$?
There are realizations by posets for $BPL(n)$ (triangulation/ball decompositions of the marked sphere) but they seem too flexible.
Maybe I need to reduce to only some triangulations that are related to the signed Schubert cell picture.
No idea how these pictures are related though.

\section{Signed Schubert Cells}

Fix $n \geq 0$. We write $\mathbf{n} = \{0, \ldots, n-1\}$.

\begin{definition}[Shape and Submatrices]
    A shape is an strictly increasing sequence of integers $0 \leq \sigma_1 < \ldots < \sigma_n$.
    A submatrix of this shape is the data of subsets $I \subseteq \mathbf{n}$ and $J \subseteq \mathbb{N}$ such that $|I| = |J|$ and $\max(J) < \sigma_{\min(I)}$.
    We denote by $\Sub(\sigma)$ the set of submatrices of shape $\sigma$.
\end{definition}

\begin{definition}[Cell]
    Let $\sigma$ be a shape and $\epsilon\in \Sub(\sigma) \to \{-, 0, +\}$ be a function.
    The signed Schubert cell $e_{\sigma, \epsilon}$ is the subset of $\Gr_n$ of all subspaces which row-reduce to $\sigma$ in the canonical basis and for which the sign of the determinant of the submatrices are given by $\epsilon$.
\end{definition}

\begin{lemma}
    All signed Schubert cells $e_{\sigma, \epsilon}$ are either empty or homeomorphic to $\mathbb{R}^m$ for some $m$.
\end{lemma}
\begin{proof}
    Idea: filration on each coefficient. Each time the fiber should be either: empty, a point or an open interval.
    This is essentially because we have a bunch of \textbf{linear} inequalities on the coefficients of the matrix, which we can rewrite as $\min < x < \max$, or equalities that have unique solutions.

    Take all submatrices involving some fix $x_{ij}$, \textit{i.e.} with $(i,j)\inset I \times J$, with a non-zero minor at this entry. We have either of the following cases:
    \begin{enumerate}
        \item one of them is sent to $0$ by the sign function $\epsilon$, in which case $x_{ij}$ is fixed to a single value, or no value.
        \item otherwise, each of them gives an open inequality on $x_{ij}$, the collection of which results in an open interval of possibilities, potentially an empty one.
    \end{enumerate} 
\end{proof}

This effectively proves that the signed Schubert cells are "elementary sets" in the sense of Aaron Peterson's 291 course, with any ordering on the coordinates.

\begin{lemma}[Faces of Schubert Cells]
    When $x_{ij}$ approaches its upper limit, it will approach a collection of inequalities. These effectively become equalities, the corresponding submatrices go from $\pm$ to $0$.
    When "$1$ approaches $0$", we have a reduction operation also giving us a single face (that's why we need all the determinants).
    Hence a "face" is choosing a subset of the $\pm$ submatrices and making them $0$.
    There's potentially $2^{\# \text{ of } \pm}$ faces, but some of them may be empty.
\end{lemma}

What is the link? Suppose $A$ is a face of $B$.
\begin{enumerate}
    \item Does all of $A$ touches $B$? Yes, do some infenitesimal perturbation of the coordinates of $A$, and you'll land in $B$.
    \item Is the link conctractible? Yes I think so.
\end{enumerate}

Then this means that the homotopy type of $BO(n)$ is the homotopy colimit of the poset of (non-empty) cells.

\section{Action on Categories with Adjoints}
\begin{enumerate}
    \item \textbf{Case $n=1$}: We have $BO(1) = B(\mathbb{Z}/2\mathbb{Z})$ which acts by taking "op" on usual $\infty$-categories.
    Say a 1-cell is sent to "op" exactly when the unique sign is negative.
    Each "2-cell" is of the form $10\ldots0\pm0\ldots0\pm0\ldots0$.
    These have 3 faces.
    Say $-+$ gives a shortened $-$, a long $0+$ and a long $-0$.
    It's easy to see there's always an even number of $-$ signs, so the action is well-defined on the 2-cells.
    It's easy to see how it extends to all cells as these 2-cells are sent to identities.

    \item \textbf{Case $n=2$}:
\end{enumerate}
\end{document}